% 論文で引用する数値. scripts/make_thesis_numbers.py で生成 (手で編集しない).

% 出典: results_gth/VALIDATION_REPORT.md の V-Sim の表 (実験 0 の CSV はなく, 結果はこの表にだけある)
% 実験 0 (ベースモデル) で DES と照合した ρ の点の数
\newcommand{\ExpZeroBasePoints}{10}

% 出典: results_gth/VALIDATION_REPORT.md の V-Sim の表 (実験 0 の CSV はなく, 結果はこの表にだけある)
% 照合した指標の数 (P_block, E[N], E[W], λ_eff, ρ_server, Cost)
\newcommand{\ExpZeroMetrics}{6}

% 出典: results_gth/VALIDATION_REPORT.md の V-Sim の表 (実験 0 の CSV はなく, 結果はこの表にだけある)
% 理論値が DES の 95% 信頼区間に入った (点, 指標) の組の数 (全点・全指標が CI 内)
\newcommand{\ExpZeroBaseInCI}{60}

% 出典: results/protect_des/protect_des.csv (policy=both)
% 6 指標の理論値が DES の 95% 信頼区間に入った組の数 (n_in_ci_6 の和)
\newcommand{\ExpZeroProtectBothInCI}{29}

% 出典: results/protect_des/protect_des.csv (policy=both)
% 照合した組の総数 (ρ の点の数 × 6 指標)
\newcommand{\ExpZeroProtectBothTotal}{30}

% 出典: results/protect_des/protect_des.csv (policy=nocancel)
% 6 指標の理論値が DES の 95% 信頼区間に入った組の数 (n_in_ci_6 の和)
\newcommand{\ExpZeroProtectNoCancelInCI}{28}

% 出典: results/protect_des/protect_des.csv (policy=nocancel)
% 照合した組の総数 (ρ の点の数 × 6 指標)
\newcommand{\ExpZeroProtectNoCancelTotal}{30}

% 出典: results/protect_des/protect_des.csv
% DES でブロックが 1 回も観測されず (CI が [0, 0]) CI 外と判定された P_block の数
\newcommand{\ExpZeroProtectUnobservable}{1}

% 出典: results/experiment_7/experiment_7_C1.csv と experiment_7_refine_C1.csv (scripts/blocking_reassessment.py の load_exp7 と assess)
% 条件 C1, 系列 Base の ERP 最小値 (Base は β を細かくした 121 点を含む)
\newcommand{\ExpSevenCOneBaseERPMin}{6.2151}

% 出典: results/experiment_7/experiment_7_C1.csv と experiment_7_refine_C1.csv (scripts/blocking_reassessment.py の load_exp7 と assess)
% 条件 C1, 系列 NoCancel の ERP 最小値 (Base は β を細かくした 121 点を含む)
\newcommand{\ExpSevenCOneNoCancelERPMin}{6.2519}

% 出典: results/experiment_7/experiment_7_C1.csv と experiment_7_refine_C1.csv (scripts/blocking_reassessment.py の load_exp7 と assess)
% 条件 C1, 系列 NoCancel の ERP 最小値の, Base の ERP 最小値に対する相対差
\newcommand{\ExpSevenCOneNoCancelERPRel}{$+0.59$\%}

% 出典: results/experiment_7/experiment_7_C1.csv と experiment_7_refine_C1.csv (scripts/blocking_reassessment.py の load_exp7 と assess)
% 条件 C1, 系列 NoCancel の ERP 最小点の P_block_arrival_stable の, Base の ERP 最小点の値に対する比
\newcommand{\ExpSevenCOneNoCancelPRatio}{1.04}

% 出典: results/experiment_7/experiment_7_C1.csv と experiment_7_refine_C1.csv (scripts/blocking_reassessment.py の load_exp7 と assess)
% 条件 C1, 系列 NoCancel の同じ信頼性での改善: P ≤ P_ref (Base の ERP 最小点の P) を
% 満たす点の中の ERP 最小値の, Base の ERP 最小値に対する −(相対差)
\newcommand{\ExpSevenCOneNoCancelSameRelImprovement}{$-0.60$\%}

% 出典: results/experiment_7/experiment_7_C1.csv と experiment_7_refine_C1.csv (scripts/blocking_reassessment.py の load_exp7 と assess)
% 条件 C1, 系列 P2 only の ERP 最小値 (Base は β を細かくした 121 点を含む)
\newcommand{\ExpSevenCOnePTwoOnlyERPMin}{6.2085}

% 出典: results/experiment_7/experiment_7_C1.csv と experiment_7_refine_C1.csv (scripts/blocking_reassessment.py の load_exp7 と assess)
% 条件 C1, 系列 P2 only の ERP 最小値の, Base の ERP 最小値に対する相対差
\newcommand{\ExpSevenCOnePTwoOnlyERPRel}{$-0.11$\%}

% 出典: results/experiment_7/experiment_7_C1.csv と experiment_7_refine_C1.csv (scripts/blocking_reassessment.py の load_exp7 と assess)
% 条件 C1, 系列 P2 only の ERP 最小点の P_block_arrival_stable の, Base の ERP 最小点の値に対する比
\newcommand{\ExpSevenCOnePTwoOnlyPRatio}{1.85}

% 出典: results/experiment_7/experiment_7_C1.csv と experiment_7_refine_C1.csv (scripts/blocking_reassessment.py の load_exp7 と assess)
% 条件 C1, 系列 P2 only の同じ信頼性での改善: P ≤ P_ref (Base の ERP 最小点の P) を
% 満たす点の中の ERP 最小値の, Base の ERP 最小値に対する −(相対差)
\newcommand{\ExpSevenCOnePTwoOnlySameRelImprovement}{$-0.01$\%}

% 出典: results/experiment_7/experiment_7_C1.csv と experiment_7_refine_C1.csv (scripts/blocking_reassessment.py の load_exp7 と assess)
% 条件 C1, 系列 P1 only の ERP 最小値 (Base は β を細かくした 121 点を含む)
\newcommand{\ExpSevenCOnePOneOnlyERPMin}{6.2066}

% 出典: results/experiment_7/experiment_7_C1.csv と experiment_7_refine_C1.csv (scripts/blocking_reassessment.py の load_exp7 と assess)
% 条件 C1, 系列 P1 only の ERP 最小値の, Base の ERP 最小値に対する相対差
\newcommand{\ExpSevenCOnePOneOnlyERPRel}{$-0.14$\%}

% 出典: results/experiment_7/experiment_7_C1.csv と experiment_7_refine_C1.csv (scripts/blocking_reassessment.py の load_exp7 と assess)
% 条件 C1, 系列 P1 only の ERP 最小点の P_block_arrival_stable の, Base の ERP 最小点の値に対する比
\newcommand{\ExpSevenCOnePOneOnlyPRatio}{0.86}

% 出典: results/experiment_7/experiment_7_C1.csv と experiment_7_refine_C1.csv (scripts/blocking_reassessment.py の load_exp7 と assess)
% 条件 C1, 系列 P1 only の同じ信頼性での改善: P ≤ P_ref (Base の ERP 最小点の P) を
% 満たす点の中の ERP 最小値の, Base の ERP 最小値に対する −(相対差)
\newcommand{\ExpSevenCOnePOneOnlySameRelImprovement}{$+0.14$\%}

% 出典: results/experiment_7/experiment_7_C1.csv と experiment_7_refine_C1.csv (scripts/blocking_reassessment.py の load_exp7 と assess)
% 条件 C1, 系列 P1 only (protect) の ERP 最小値 (Base は β を細かくした 121 点を含む)
\newcommand{\ExpSevenCOnePOneOnlyProtectERPMin}{6.1432}

% 出典: results/experiment_7/experiment_7_C1.csv と experiment_7_refine_C1.csv (scripts/blocking_reassessment.py の load_exp7 と assess)
% 条件 C1, 系列 P1 only (protect) の ERP 最小値の, Base の ERP 最小値に対する相対差
\newcommand{\ExpSevenCOnePOneOnlyProtectERPRel}{$-1.16$\%}

% 出典: results/experiment_7/experiment_7_C1.csv と experiment_7_refine_C1.csv (scripts/blocking_reassessment.py の load_exp7 と assess)
% 条件 C1, 系列 P1 only (protect) の ERP 最小点の P_block_arrival_stable の, Base の ERP 最小点の値に対する比
\newcommand{\ExpSevenCOnePOneOnlyProtectPRatio}{2.35}

% 出典: results/experiment_7/experiment_7_C1.csv と experiment_7_refine_C1.csv (scripts/blocking_reassessment.py の load_exp7 と assess)
% 条件 C1, 系列 P1 only (protect) の同じ信頼性での改善: P ≤ P_ref (Base の ERP 最小点の P) を
% 満たす点の中の ERP 最小値の, Base の ERP 最小値に対する −(相対差)
\newcommand{\ExpSevenCOnePOneOnlyProtectSameRelImprovement}{$+0.11$\%}

% 出典: results/experiment_7/experiment_7_C1.csv と experiment_7_refine_C1.csv (scripts/blocking_reassessment.py の load_exp7 と assess)
% 条件 C1, 系列 P1+P2 の ERP 最小値 (Base は β を細かくした 121 点を含む)
\newcommand{\ExpSevenCOnePOneTwoERPMin}{6.1936}

% 出典: results/experiment_7/experiment_7_C1.csv と experiment_7_refine_C1.csv (scripts/blocking_reassessment.py の load_exp7 と assess)
% 条件 C1, 系列 P1+P2 の ERP 最小値の, Base の ERP 最小値に対する相対差
\newcommand{\ExpSevenCOnePOneTwoERPRel}{$-0.35$\%}

% 出典: results/experiment_7/experiment_7_C1.csv と experiment_7_refine_C1.csv (scripts/blocking_reassessment.py の load_exp7 と assess)
% 条件 C1, 系列 P1+P2 の ERP 最小点の P_block_arrival_stable の, Base の ERP 最小点の値に対する比
\newcommand{\ExpSevenCOnePOneTwoPRatio}{1.98}

% 出典: results/experiment_7/experiment_7_C1.csv と experiment_7_refine_C1.csv (scripts/blocking_reassessment.py の load_exp7 と assess)
% 条件 C1, 系列 P1+P2 の同じ信頼性での改善: P ≤ P_ref (Base の ERP 最小点の P) を
% 満たす点の中の ERP 最小値の, Base の ERP 最小値に対する −(相対差)
\newcommand{\ExpSevenCOnePOneTwoSameRelImprovement}{$+0.08$\%}

% 出典: results/experiment_7/experiment_7_C1.csv と experiment_7_refine_C1.csv (scripts/blocking_reassessment.py の load_exp7 と assess)
% 条件 C1, 系列 P1+P2 (protect) の ERP 最小値 (Base は β を細かくした 121 点を含む)
\newcommand{\ExpSevenCOnePOneTwoProtectERPMin}{6.1409}

% 出典: results/experiment_7/experiment_7_C1.csv と experiment_7_refine_C1.csv (scripts/blocking_reassessment.py の load_exp7 と assess)
% 条件 C1, 系列 P1+P2 (protect) の ERP 最小値の, Base の ERP 最小値に対する相対差
\newcommand{\ExpSevenCOnePOneTwoProtectERPRel}{$-1.19$\%}

% 出典: results/experiment_7/experiment_7_C1.csv と experiment_7_refine_C1.csv (scripts/blocking_reassessment.py の load_exp7 と assess)
% 条件 C1, 系列 P1+P2 (protect) の ERP 最小点の P_block_arrival_stable の, Base の ERP 最小点の値に対する比
\newcommand{\ExpSevenCOnePOneTwoProtectPRatio}{3.07}

% 出典: results/experiment_7/experiment_7_C1.csv と experiment_7_refine_C1.csv (scripts/blocking_reassessment.py の load_exp7 と assess)
% 条件 C1, 系列 P1+P2 (protect) の同じ信頼性での改善: P ≤ P_ref (Base の ERP 最小点の P) を
% 満たす点の中の ERP 最小値の, Base の ERP 最小値に対する −(相対差)
\newcommand{\ExpSevenCOnePOneTwoProtectSameRelImprovement}{$+0.81$\%}

% 出典: results/experiment_7/experiment_7_C1.csv と experiment_7_refine_C1.csv (scripts/blocking_reassessment.py の load_exp7 と assess)
% 条件 C1, Base の ERP 最小点の P_block_arrival_stable (= P_ref)
\newcommand{\ExpSevenCOneBasePBlock}{0.000656}

% 出典: results/experiment_7/experiment_7_C1.csv と experiment_7_refine_C1.csv (scripts/blocking_reassessment.py の load_exp7 と assess)
% 条件 C1, P_block_arrival_stable ≤ 1e-2 を満たす点をもつ系列の数 (Base を含む 7 系列中)
\newcommand{\ExpSevenCOneEpsTwoFeasibleCount}{7}

% 出典: results/experiment_7/experiment_7_C2.csv と experiment_7_refine_C2.csv (scripts/blocking_reassessment.py の load_exp7 と assess)
% 条件 C2, 系列 Base の ERP 最小値 (Base は β を細かくした 121 点を含む)
\newcommand{\ExpSevenCTwoBaseERPMin}{9.5494}

% 出典: results/experiment_7/experiment_7_C2.csv と experiment_7_refine_C2.csv (scripts/blocking_reassessment.py の load_exp7 と assess)
% 条件 C2, 系列 NoCancel の ERP 最小値 (Base は β を細かくした 121 点を含む)
\newcommand{\ExpSevenCTwoNoCancelERPMin}{9.7522}

% 出典: results/experiment_7/experiment_7_C2.csv と experiment_7_refine_C2.csv (scripts/blocking_reassessment.py の load_exp7 と assess)
% 条件 C2, 系列 NoCancel の ERP 最小値の, Base の ERP 最小値に対する相対差
\newcommand{\ExpSevenCTwoNoCancelERPRel}{$+2.12$\%}

% 出典: results/experiment_7/experiment_7_C2.csv と experiment_7_refine_C2.csv (scripts/blocking_reassessment.py の load_exp7 と assess)
% 条件 C2, 系列 NoCancel の ERP 最小点の P_block_arrival_stable の, Base の ERP 最小点の値に対する比
\newcommand{\ExpSevenCTwoNoCancelPRatio}{0.92}

% 出典: results/experiment_7/experiment_7_C2.csv と experiment_7_refine_C2.csv (scripts/blocking_reassessment.py の load_exp7 と assess)
% 条件 C2, 系列 NoCancel の同じ信頼性での改善: P ≤ P_ref (Base の ERP 最小点の P) を
% 満たす点の中の ERP 最小値の, Base の ERP 最小値に対する −(相対差)
\newcommand{\ExpSevenCTwoNoCancelSameRelImprovement}{$-2.12$\%}

% 出典: results/experiment_7/experiment_7_C2.csv と experiment_7_refine_C2.csv (scripts/blocking_reassessment.py の load_exp7 と assess)
% 条件 C2, 系列 P2 only の ERP 最小値 (Base は β を細かくした 121 点を含む)
\newcommand{\ExpSevenCTwoPTwoOnlyERPMin}{9.5384}

% 出典: results/experiment_7/experiment_7_C2.csv と experiment_7_refine_C2.csv (scripts/blocking_reassessment.py の load_exp7 と assess)
% 条件 C2, 系列 P2 only の ERP 最小値の, Base の ERP 最小値に対する相対差
\newcommand{\ExpSevenCTwoPTwoOnlyERPRel}{$-0.11$\%}

% 出典: results/experiment_7/experiment_7_C2.csv と experiment_7_refine_C2.csv (scripts/blocking_reassessment.py の load_exp7 と assess)
% 条件 C2, 系列 P2 only の ERP 最小点の P_block_arrival_stable の, Base の ERP 最小点の値に対する比
\newcommand{\ExpSevenCTwoPTwoOnlyPRatio}{1.01}

% 出典: results/experiment_7/experiment_7_C2.csv と experiment_7_refine_C2.csv (scripts/blocking_reassessment.py の load_exp7 と assess)
% 条件 C2, 系列 P2 only の同じ信頼性での改善: P ≤ P_ref (Base の ERP 最小点の P) を
% 満たす点の中の ERP 最小値の, Base の ERP 最小値に対する −(相対差)
\newcommand{\ExpSevenCTwoPTwoOnlySameRelImprovement}{$+0.10$\%}

% 出典: results/experiment_7/experiment_7_C2.csv と experiment_7_refine_C2.csv (scripts/blocking_reassessment.py の load_exp7 と assess)
% 条件 C2, 系列 P1 only の ERP 最小値 (Base は β を細かくした 121 点を含む)
\newcommand{\ExpSevenCTwoPOneOnlyERPMin}{9.5414}

% 出典: results/experiment_7/experiment_7_C2.csv と experiment_7_refine_C2.csv (scripts/blocking_reassessment.py の load_exp7 と assess)
% 条件 C2, 系列 P1 only の ERP 最小値の, Base の ERP 最小値に対する相対差
\newcommand{\ExpSevenCTwoPOneOnlyERPRel}{$-0.08$\%}

% 出典: results/experiment_7/experiment_7_C2.csv と experiment_7_refine_C2.csv (scripts/blocking_reassessment.py の load_exp7 と assess)
% 条件 C2, 系列 P1 only の ERP 最小点の P_block_arrival_stable の, Base の ERP 最小点の値に対する比
\newcommand{\ExpSevenCTwoPOneOnlyPRatio}{0.98}

% 出典: results/experiment_7/experiment_7_C2.csv と experiment_7_refine_C2.csv (scripts/blocking_reassessment.py の load_exp7 と assess)
% 条件 C2, 系列 P1 only の同じ信頼性での改善: P ≤ P_ref (Base の ERP 最小点の P) を
% 満たす点の中の ERP 最小値の, Base の ERP 最小値に対する −(相対差)
\newcommand{\ExpSevenCTwoPOneOnlySameRelImprovement}{$+0.08$\%}

% 出典: results/experiment_7/experiment_7_C2.csv と experiment_7_refine_C2.csv (scripts/blocking_reassessment.py の load_exp7 と assess)
% 条件 C2, 系列 P1 only (protect) の ERP 最小値 (Base は β を細かくした 121 点を含む)
\newcommand{\ExpSevenCTwoPOneOnlyProtectERPMin}{9.5193}

% 出典: results/experiment_7/experiment_7_C2.csv と experiment_7_refine_C2.csv (scripts/blocking_reassessment.py の load_exp7 と assess)
% 条件 C2, 系列 P1 only (protect) の ERP 最小値の, Base の ERP 最小値に対する相対差
\newcommand{\ExpSevenCTwoPOneOnlyProtectERPRel}{$-0.32$\%}

% 出典: results/experiment_7/experiment_7_C2.csv と experiment_7_refine_C2.csv (scripts/blocking_reassessment.py の load_exp7 と assess)
% 条件 C2, 系列 P1 only (protect) の ERP 最小点の P_block_arrival_stable の, Base の ERP 最小点の値に対する比
\newcommand{\ExpSevenCTwoPOneOnlyProtectPRatio}{0.96}

% 出典: results/experiment_7/experiment_7_C2.csv と experiment_7_refine_C2.csv (scripts/blocking_reassessment.py の load_exp7 と assess)
% 条件 C2, 系列 P1 only (protect) の同じ信頼性での改善: P ≤ P_ref (Base の ERP 最小点の P) を
% 満たす点の中の ERP 最小値の, Base の ERP 最小値に対する −(相対差)
\newcommand{\ExpSevenCTwoPOneOnlyProtectSameRelImprovement}{$+0.32$\%}

% 出典: results/experiment_7/experiment_7_C2.csv と experiment_7_refine_C2.csv (scripts/blocking_reassessment.py の load_exp7 と assess)
% 条件 C2, 系列 P1+P2 の ERP 最小値 (Base は β を細かくした 121 点を含む)
\newcommand{\ExpSevenCTwoPOneTwoERPMin}{9.5293}

% 出典: results/experiment_7/experiment_7_C2.csv と experiment_7_refine_C2.csv (scripts/blocking_reassessment.py の load_exp7 と assess)
% 条件 C2, 系列 P1+P2 の ERP 最小値の, Base の ERP 最小値に対する相対差
\newcommand{\ExpSevenCTwoPOneTwoERPRel}{$-0.21$\%}

% 出典: results/experiment_7/experiment_7_C2.csv と experiment_7_refine_C2.csv (scripts/blocking_reassessment.py の load_exp7 と assess)
% 条件 C2, 系列 P1+P2 の ERP 最小点の P_block_arrival_stable の, Base の ERP 最小点の値に対する比
\newcommand{\ExpSevenCTwoPOneTwoPRatio}{0.99}

% 出典: results/experiment_7/experiment_7_C2.csv と experiment_7_refine_C2.csv (scripts/blocking_reassessment.py の load_exp7 と assess)
% 条件 C2, 系列 P1+P2 の同じ信頼性での改善: P ≤ P_ref (Base の ERP 最小点の P) を
% 満たす点の中の ERP 最小値の, Base の ERP 最小値に対する −(相対差)
\newcommand{\ExpSevenCTwoPOneTwoSameRelImprovement}{$+0.21$\%}

% 出典: results/experiment_7/experiment_7_C2.csv と experiment_7_refine_C2.csv (scripts/blocking_reassessment.py の load_exp7 と assess)
% 条件 C2, 系列 P1+P2 (protect) の ERP 最小値 (Base は β を細かくした 121 点を含む)
\newcommand{\ExpSevenCTwoPOneTwoProtectERPMin}{9.5187}

% 出典: results/experiment_7/experiment_7_C2.csv と experiment_7_refine_C2.csv (scripts/blocking_reassessment.py の load_exp7 と assess)
% 条件 C2, 系列 P1+P2 (protect) の ERP 最小値の, Base の ERP 最小値に対する相対差
\newcommand{\ExpSevenCTwoPOneTwoProtectERPRel}{$-0.32$\%}

% 出典: results/experiment_7/experiment_7_C2.csv と experiment_7_refine_C2.csv (scripts/blocking_reassessment.py の load_exp7 と assess)
% 条件 C2, 系列 P1+P2 (protect) の ERP 最小点の P_block_arrival_stable の, Base の ERP 最小点の値に対する比
\newcommand{\ExpSevenCTwoPOneTwoProtectPRatio}{0.97}

% 出典: results/experiment_7/experiment_7_C2.csv と experiment_7_refine_C2.csv (scripts/blocking_reassessment.py の load_exp7 と assess)
% 条件 C2, 系列 P1+P2 (protect) の同じ信頼性での改善: P ≤ P_ref (Base の ERP 最小点の P) を
% 満たす点の中の ERP 最小値の, Base の ERP 最小値に対する −(相対差)
\newcommand{\ExpSevenCTwoPOneTwoProtectSameRelImprovement}{$+0.32$\%}

% 出典: results/experiment_7/experiment_7_C2.csv と experiment_7_refine_C2.csv (scripts/blocking_reassessment.py の load_exp7 と assess)
% 条件 C2, Base の ERP 最小点の P_block_arrival_stable (= P_ref)
\newcommand{\ExpSevenCTwoBasePBlock}{0.129}

% 出典: results/experiment_7/experiment_7_C2.csv と experiment_7_refine_C2.csv (scripts/blocking_reassessment.py の load_exp7 と assess)
% 条件 C2, P_block_arrival_stable ≤ 1e-2 を満たす点をもつ系列の数 (Base を含む 7 系列中)
\newcommand{\ExpSevenCTwoEpsTwoFeasibleCount}{0}

% 出典: results/experiment_7/experiment_7_C3.csv と experiment_7_refine_C3.csv (scripts/blocking_reassessment.py の load_exp7 と assess)
% 条件 C3, 系列 Base の ERP 最小値 (Base は β を細かくした 121 点を含む)
\newcommand{\ExpSevenCThreeBaseERPMin}{5.6791}

% 出典: results/experiment_7/experiment_7_C3.csv と experiment_7_refine_C3.csv (scripts/blocking_reassessment.py の load_exp7 と assess)
% 条件 C3, 系列 NoCancel の ERP 最小値 (Base は β を細かくした 121 点を含む)
\newcommand{\ExpSevenCThreeNoCancelERPMin}{6.1654}

% 出典: results/experiment_7/experiment_7_C3.csv と experiment_7_refine_C3.csv (scripts/blocking_reassessment.py の load_exp7 と assess)
% 条件 C3, 系列 NoCancel の ERP 最小値の, Base の ERP 最小値に対する相対差
\newcommand{\ExpSevenCThreeNoCancelERPRel}{$+8.56$\%}

% 出典: results/experiment_7/experiment_7_C3.csv と experiment_7_refine_C3.csv (scripts/blocking_reassessment.py の load_exp7 と assess)
% 条件 C3, 系列 NoCancel の ERP 最小点の P_block_arrival_stable の, Base の ERP 最小点の値に対する比
\newcommand{\ExpSevenCThreeNoCancelPRatio}{0.08}

% 出典: results/experiment_7/experiment_7_C3.csv と experiment_7_refine_C3.csv (scripts/blocking_reassessment.py の load_exp7 と assess)
% 条件 C3, 系列 NoCancel の同じ信頼性での改善: P ≤ P_ref (Base の ERP 最小点の P) を
% 満たす点の中の ERP 最小値の, Base の ERP 最小値に対する −(相対差)
\newcommand{\ExpSevenCThreeNoCancelSameRelImprovement}{$-8.56$\%}

% 出典: results/experiment_7/experiment_7_C3.csv と experiment_7_refine_C3.csv (scripts/blocking_reassessment.py の load_exp7 と assess)
% 条件 C3, 系列 P2 only の ERP 最小値 (Base は β を細かくした 121 点を含む)
\newcommand{\ExpSevenCThreePTwoOnlyERPMin}{5.6800}

% 出典: results/experiment_7/experiment_7_C3.csv と experiment_7_refine_C3.csv (scripts/blocking_reassessment.py の load_exp7 と assess)
% 条件 C3, 系列 P2 only の ERP 最小値の, Base の ERP 最小値に対する相対差
\newcommand{\ExpSevenCThreePTwoOnlyERPRel}{$+0.02$\%}

% 出典: results/experiment_7/experiment_7_C3.csv と experiment_7_refine_C3.csv (scripts/blocking_reassessment.py の load_exp7 と assess)
% 条件 C3, 系列 P2 only の ERP 最小点の P_block_arrival_stable の, Base の ERP 最小点の値に対する比
\newcommand{\ExpSevenCThreePTwoOnlyPRatio}{1.15}

% 出典: results/experiment_7/experiment_7_C3.csv と experiment_7_refine_C3.csv (scripts/blocking_reassessment.py の load_exp7 と assess)
% 条件 C3, 系列 P2 only の同じ信頼性での改善: P ≤ P_ref (Base の ERP 最小点の P) を
% 満たす点の中の ERP 最小値の, Base の ERP 最小値に対する −(相対差)
\newcommand{\ExpSevenCThreePTwoOnlySameRelImprovement}{$-1.21$\%}

% 出典: results/experiment_7/experiment_7_C3.csv と experiment_7_refine_C3.csv (scripts/blocking_reassessment.py の load_exp7 と assess)
% 条件 C3, 系列 P1 only の ERP 最小値 (Base は β を細かくした 121 点を含む)
\newcommand{\ExpSevenCThreePOneOnlyERPMin}{5.6795}

% 出典: results/experiment_7/experiment_7_C3.csv と experiment_7_refine_C3.csv (scripts/blocking_reassessment.py の load_exp7 と assess)
% 条件 C3, 系列 P1 only の ERP 最小値の, Base の ERP 最小値に対する相対差
\newcommand{\ExpSevenCThreePOneOnlyERPRel}{$+0.01$\%}

% 出典: results/experiment_7/experiment_7_C3.csv と experiment_7_refine_C3.csv (scripts/blocking_reassessment.py の load_exp7 と assess)
% 条件 C3, 系列 P1 only の ERP 最小点の P_block_arrival_stable の, Base の ERP 最小点の値に対する比
\newcommand{\ExpSevenCThreePOneOnlyPRatio}{1.18}

% 出典: results/experiment_7/experiment_7_C3.csv と experiment_7_refine_C3.csv (scripts/blocking_reassessment.py の load_exp7 と assess)
% 条件 C3, 系列 P1 only の同じ信頼性での改善: P ≤ P_ref (Base の ERP 最小点の P) を
% 満たす点の中の ERP 最小値の, Base の ERP 最小値に対する −(相対差)
\newcommand{\ExpSevenCThreePOneOnlySameRelImprovement}{$-0.90$\%}

% 出典: results/experiment_7/experiment_7_C3.csv と experiment_7_refine_C3.csv (scripts/blocking_reassessment.py の load_exp7 と assess)
% 条件 C3, 系列 P1 only (protect) の ERP 最小値 (Base は β を細かくした 121 点を含む)
\newcommand{\ExpSevenCThreePOneOnlyProtectERPMin}{5.6632}

% 出典: results/experiment_7/experiment_7_C3.csv と experiment_7_refine_C3.csv (scripts/blocking_reassessment.py の load_exp7 と assess)
% 条件 C3, 系列 P1 only (protect) の ERP 最小値の, Base の ERP 最小値に対する相対差
\newcommand{\ExpSevenCThreePOneOnlyProtectERPRel}{$-0.28$\%}

% 出典: results/experiment_7/experiment_7_C3.csv と experiment_7_refine_C3.csv (scripts/blocking_reassessment.py の load_exp7 と assess)
% 条件 C3, 系列 P1 only (protect) の ERP 最小点の P_block_arrival_stable の, Base の ERP 最小点の値に対する比
\newcommand{\ExpSevenCThreePOneOnlyProtectPRatio}{0.69}

% 出典: results/experiment_7/experiment_7_C3.csv と experiment_7_refine_C3.csv (scripts/blocking_reassessment.py の load_exp7 と assess)
% 条件 C3, 系列 P1 only (protect) の同じ信頼性での改善: P ≤ P_ref (Base の ERP 最小点の P) を
% 満たす点の中の ERP 最小値の, Base の ERP 最小値に対する −(相対差)
\newcommand{\ExpSevenCThreePOneOnlyProtectSameRelImprovement}{$+0.28$\%}

% 出典: results/experiment_7/experiment_7_C3.csv と experiment_7_refine_C3.csv (scripts/blocking_reassessment.py の load_exp7 と assess)
% 条件 C3, 系列 P1+P2 の ERP 最小値 (Base は β を細かくした 121 点を含む)
\newcommand{\ExpSevenCThreePOneTwoERPMin}{5.6789}

% 出典: results/experiment_7/experiment_7_C3.csv と experiment_7_refine_C3.csv (scripts/blocking_reassessment.py の load_exp7 と assess)
% 条件 C3, 系列 P1+P2 の ERP 最小値の, Base の ERP 最小値に対する相対差
\newcommand{\ExpSevenCThreePOneTwoERPRel}{$-0.00$\%}

% 出典: results/experiment_7/experiment_7_C3.csv と experiment_7_refine_C3.csv (scripts/blocking_reassessment.py の load_exp7 と assess)
% 条件 C3, 系列 P1+P2 の ERP 最小点の P_block_arrival_stable の, Base の ERP 最小点の値に対する比
\newcommand{\ExpSevenCThreePOneTwoPRatio}{1.00}

% 出典: results/experiment_7/experiment_7_C3.csv と experiment_7_refine_C3.csv (scripts/blocking_reassessment.py の load_exp7 と assess)
% 条件 C3, 系列 P1+P2 の同じ信頼性での改善: P ≤ P_ref (Base の ERP 最小点の P) を
% 満たす点の中の ERP 最小値の, Base の ERP 最小値に対する −(相対差)
\newcommand{\ExpSevenCThreePOneTwoSameRelImprovement}{$-1.21$\%}

% 出典: results/experiment_7/experiment_7_C3.csv と experiment_7_refine_C3.csv (scripts/blocking_reassessment.py の load_exp7 と assess)
% 条件 C3, 系列 P1+P2 (protect) の ERP 最小値 (Base は β を細かくした 121 点を含む)
\newcommand{\ExpSevenCThreePOneTwoProtectERPMin}{5.6736}

% 出典: results/experiment_7/experiment_7_C3.csv と experiment_7_refine_C3.csv (scripts/blocking_reassessment.py の load_exp7 と assess)
% 条件 C3, 系列 P1+P2 (protect) の ERP 最小値の, Base の ERP 最小値に対する相対差
\newcommand{\ExpSevenCThreePOneTwoProtectERPRel}{$-0.10$\%}

% 出典: results/experiment_7/experiment_7_C3.csv と experiment_7_refine_C3.csv (scripts/blocking_reassessment.py の load_exp7 と assess)
% 条件 C3, 系列 P1+P2 (protect) の ERP 最小点の P_block_arrival_stable の, Base の ERP 最小点の値に対する比
\newcommand{\ExpSevenCThreePOneTwoProtectPRatio}{0.68}

% 出典: results/experiment_7/experiment_7_C3.csv と experiment_7_refine_C3.csv (scripts/blocking_reassessment.py の load_exp7 と assess)
% 条件 C3, 系列 P1+P2 (protect) の同じ信頼性での改善: P ≤ P_ref (Base の ERP 最小点の P) を
% 満たす点の中の ERP 最小値の, Base の ERP 最小値に対する −(相対差)
\newcommand{\ExpSevenCThreePOneTwoProtectSameRelImprovement}{$+0.10$\%}

% 出典: results/experiment_7/experiment_7_C3.csv と experiment_7_refine_C3.csv (scripts/blocking_reassessment.py の load_exp7 と assess)
% 条件 C3, Base の ERP 最小点の P_block_arrival_stable (= P_ref)
\newcommand{\ExpSevenCThreeBasePBlock}{6.68e-05}

% 出典: results/experiment_7/experiment_7_C3.csv と experiment_7_refine_C3.csv (scripts/blocking_reassessment.py の load_exp7 と assess)
% 条件 C3, P_block_arrival_stable ≤ 1e-2 を満たす点をもつ系列の数 (Base を含む 7 系列中)
\newcommand{\ExpSevenCThreeEpsTwoFeasibleCount}{7}

% 出典: results/experiment_7/experiment_7_C4.csv と experiment_7_refine_C4.csv (scripts/blocking_reassessment.py の load_exp7 と assess)
% 条件 C4, 系列 Base の ERP 最小値 (Base は β を細かくした 121 点を含む)
\newcommand{\ExpSevenCFourBaseERPMin}{6.2891}

% 出典: results/experiment_7/experiment_7_C4.csv と experiment_7_refine_C4.csv (scripts/blocking_reassessment.py の load_exp7 と assess)
% 条件 C4, 系列 NoCancel の ERP 最小値 (Base は β を細かくした 121 点を含む)
\newcommand{\ExpSevenCFourNoCancelERPMin}{6.3099}

% 出典: results/experiment_7/experiment_7_C4.csv と experiment_7_refine_C4.csv (scripts/blocking_reassessment.py の load_exp7 と assess)
% 条件 C4, 系列 NoCancel の ERP 最小値の, Base の ERP 最小値に対する相対差
\newcommand{\ExpSevenCFourNoCancelERPRel}{$+0.33$\%}

% 出典: results/experiment_7/experiment_7_C4.csv と experiment_7_refine_C4.csv (scripts/blocking_reassessment.py の load_exp7 と assess)
% 条件 C4, 系列 NoCancel の ERP 最小点の P_block_arrival_stable の, Base の ERP 最小点の値に対する比
\newcommand{\ExpSevenCFourNoCancelPRatio}{0.95}

% 出典: results/experiment_7/experiment_7_C4.csv と experiment_7_refine_C4.csv (scripts/blocking_reassessment.py の load_exp7 と assess)
% 条件 C4, 系列 NoCancel の同じ信頼性での改善: P ≤ P_ref (Base の ERP 最小点の P) を
% 満たす点の中の ERP 最小値の, Base の ERP 最小値に対する −(相対差)
\newcommand{\ExpSevenCFourNoCancelSameRelImprovement}{$-0.33$\%}

% 出典: results/experiment_7/experiment_7_C4.csv と experiment_7_refine_C4.csv (scripts/blocking_reassessment.py の load_exp7 と assess)
% 条件 C4, 系列 P2 only の ERP 最小値 (Base は β を細かくした 121 点を含む)
\newcommand{\ExpSevenCFourPTwoOnlyERPMin}{6.2826}

% 出典: results/experiment_7/experiment_7_C4.csv と experiment_7_refine_C4.csv (scripts/blocking_reassessment.py の load_exp7 と assess)
% 条件 C4, 系列 P2 only の ERP 最小値の, Base の ERP 最小値に対する相対差
\newcommand{\ExpSevenCFourPTwoOnlyERPRel}{$-0.10$\%}

% 出典: results/experiment_7/experiment_7_C4.csv と experiment_7_refine_C4.csv (scripts/blocking_reassessment.py の load_exp7 と assess)
% 条件 C4, 系列 P2 only の ERP 最小点の P_block_arrival_stable の, Base の ERP 最小点の値に対する比
\newcommand{\ExpSevenCFourPTwoOnlyPRatio}{1.88}

% 出典: results/experiment_7/experiment_7_C4.csv と experiment_7_refine_C4.csv (scripts/blocking_reassessment.py の load_exp7 と assess)
% 条件 C4, 系列 P2 only の同じ信頼性での改善: P ≤ P_ref (Base の ERP 最小点の P) を
% 満たす点の中の ERP 最小値の, Base の ERP 最小値に対する −(相対差)
\newcommand{\ExpSevenCFourPTwoOnlySameRelImprovement}{$+0.07$\%}

% 出典: results/experiment_7/experiment_7_C4.csv と experiment_7_refine_C4.csv (scripts/blocking_reassessment.py の load_exp7 と assess)
% 条件 C4, 系列 P1 only の ERP 最小値 (Base は β を細かくした 121 点を含む)
\newcommand{\ExpSevenCFourPOneOnlyERPMin}{6.2938}

% 出典: results/experiment_7/experiment_7_C4.csv と experiment_7_refine_C4.csv (scripts/blocking_reassessment.py の load_exp7 と assess)
% 条件 C4, 系列 P1 only の ERP 最小値の, Base の ERP 最小値に対する相対差
\newcommand{\ExpSevenCFourPOneOnlyERPRel}{$+0.07$\%}

% 出典: results/experiment_7/experiment_7_C4.csv と experiment_7_refine_C4.csv (scripts/blocking_reassessment.py の load_exp7 と assess)
% 条件 C4, 系列 P1 only の ERP 最小点の P_block_arrival_stable の, Base の ERP 最小点の値に対する比
\newcommand{\ExpSevenCFourPOneOnlyPRatio}{0.47}

% 出典: results/experiment_7/experiment_7_C4.csv と experiment_7_refine_C4.csv (scripts/blocking_reassessment.py の load_exp7 と assess)
% 条件 C4, 系列 P1 only の同じ信頼性での改善: P ≤ P_ref (Base の ERP 最小点の P) を
% 満たす点の中の ERP 最小値の, Base の ERP 最小値に対する −(相対差)
\newcommand{\ExpSevenCFourPOneOnlySameRelImprovement}{$-0.07$\%}

% 出典: results/experiment_7/experiment_7_C4.csv と experiment_7_refine_C4.csv (scripts/blocking_reassessment.py の load_exp7 と assess)
% 条件 C4, 系列 P1 only (protect) の ERP 最小値 (Base は β を細かくした 121 点を含む)
\newcommand{\ExpSevenCFourPOneOnlyProtectERPMin}{6.2757}

% 出典: results/experiment_7/experiment_7_C4.csv と experiment_7_refine_C4.csv (scripts/blocking_reassessment.py の load_exp7 と assess)
% 条件 C4, 系列 P1 only (protect) の ERP 最小値の, Base の ERP 最小値に対する相対差
\newcommand{\ExpSevenCFourPOneOnlyProtectERPRel}{$-0.21$\%}

% 出典: results/experiment_7/experiment_7_C4.csv と experiment_7_refine_C4.csv (scripts/blocking_reassessment.py の load_exp7 と assess)
% 条件 C4, 系列 P1 only (protect) の ERP 最小点の P_block_arrival_stable の, Base の ERP 最小点の値に対する比
\newcommand{\ExpSevenCFourPOneOnlyProtectPRatio}{1.84}

% 出典: results/experiment_7/experiment_7_C4.csv と experiment_7_refine_C4.csv (scripts/blocking_reassessment.py の load_exp7 と assess)
% 条件 C4, 系列 P1 only (protect) の同じ信頼性での改善: P ≤ P_ref (Base の ERP 最小点の P) を
% 満たす点の中の ERP 最小値の, Base の ERP 最小値に対する −(相対差)
\newcommand{\ExpSevenCFourPOneOnlyProtectSameRelImprovement}{$-0.08$\%}

% 出典: results/experiment_7/experiment_7_C4.csv と experiment_7_refine_C4.csv (scripts/blocking_reassessment.py の load_exp7 と assess)
% 条件 C4, 系列 P1+P2 の ERP 最小値 (Base は β を細かくした 121 点を含む)
\newcommand{\ExpSevenCFourPOneTwoERPMin}{6.2822}

% 出典: results/experiment_7/experiment_7_C4.csv と experiment_7_refine_C4.csv (scripts/blocking_reassessment.py の load_exp7 と assess)
% 条件 C4, 系列 P1+P2 の ERP 最小値の, Base の ERP 最小値に対する相対差
\newcommand{\ExpSevenCFourPOneTwoERPRel}{$-0.11$\%}

% 出典: results/experiment_7/experiment_7_C4.csv と experiment_7_refine_C4.csv (scripts/blocking_reassessment.py の load_exp7 と assess)
% 条件 C4, 系列 P1+P2 の ERP 最小点の P_block_arrival_stable の, Base の ERP 最小点の値に対する比
\newcommand{\ExpSevenCFourPOneTwoPRatio}{1.87}

% 出典: results/experiment_7/experiment_7_C4.csv と experiment_7_refine_C4.csv (scripts/blocking_reassessment.py の load_exp7 と assess)
% 条件 C4, 系列 P1+P2 の同じ信頼性での改善: P ≤ P_ref (Base の ERP 最小点の P) を
% 満たす点の中の ERP 最小値の, Base の ERP 最小値に対する −(相対差)
\newcommand{\ExpSevenCFourPOneTwoSameRelImprovement}{$+0.07$\%}

% 出典: results/experiment_7/experiment_7_C4.csv と experiment_7_refine_C4.csv (scripts/blocking_reassessment.py の load_exp7 と assess)
% 条件 C4, 系列 P1+P2 (protect) の ERP 最小値 (Base は β を細かくした 121 点を含む)
\newcommand{\ExpSevenCFourPOneTwoProtectERPMin}{6.2749}

% 出典: results/experiment_7/experiment_7_C4.csv と experiment_7_refine_C4.csv (scripts/blocking_reassessment.py の load_exp7 と assess)
% 条件 C4, 系列 P1+P2 (protect) の ERP 最小値の, Base の ERP 最小値に対する相対差
\newcommand{\ExpSevenCFourPOneTwoProtectERPRel}{$-0.23$\%}

% 出典: results/experiment_7/experiment_7_C4.csv と experiment_7_refine_C4.csv (scripts/blocking_reassessment.py の load_exp7 と assess)
% 条件 C4, 系列 P1+P2 (protect) の ERP 最小点の P_block_arrival_stable の, Base の ERP 最小点の値に対する比
\newcommand{\ExpSevenCFourPOneTwoProtectPRatio}{2.48}

% 出典: results/experiment_7/experiment_7_C4.csv と experiment_7_refine_C4.csv (scripts/blocking_reassessment.py の load_exp7 と assess)
% 条件 C4, 系列 P1+P2 (protect) の同じ信頼性での改善: P ≤ P_ref (Base の ERP 最小点の P) を
% 満たす点の中の ERP 最小値の, Base の ERP 最小値に対する −(相対差)
\newcommand{\ExpSevenCFourPOneTwoProtectSameRelImprovement}{$+0.07$\%}

% 出典: results/experiment_7/experiment_7_C4.csv と experiment_7_refine_C4.csv (scripts/blocking_reassessment.py の load_exp7 と assess)
% 条件 C4, Base の ERP 最小点の P_block_arrival_stable (= P_ref)
\newcommand{\ExpSevenCFourBasePBlock}{0.000553}

% 出典: results/experiment_7/experiment_7_C4.csv と experiment_7_refine_C4.csv (scripts/blocking_reassessment.py の load_exp7 と assess)
% 条件 C4, P_block_arrival_stable ≤ 1e-2 を満たす点をもつ系列の数 (Base を含む 7 系列中)
\newcommand{\ExpSevenCFourEpsTwoFeasibleCount}{7}

% 出典: results/experiment_7/experiment_7_C5.csv と experiment_7_refine_C5.csv (scripts/blocking_reassessment.py の load_exp7 と assess)
% 条件 C5, 系列 Base の ERP 最小値 (Base は β を細かくした 121 点を含む)
\newcommand{\ExpSevenCFiveBaseERPMin}{4.3070}

% 出典: results/experiment_7/experiment_7_C5.csv と experiment_7_refine_C5.csv (scripts/blocking_reassessment.py の load_exp7 と assess)
% 条件 C5, 系列 NoCancel の ERP 最小値 (Base は β を細かくした 121 点を含む)
\newcommand{\ExpSevenCFiveNoCancelERPMin}{4.3498}

% 出典: results/experiment_7/experiment_7_C5.csv と experiment_7_refine_C5.csv (scripts/blocking_reassessment.py の load_exp7 と assess)
% 条件 C5, 系列 NoCancel の ERP 最小値の, Base の ERP 最小値に対する相対差
\newcommand{\ExpSevenCFiveNoCancelERPRel}{$+0.99$\%}

% 出典: results/experiment_7/experiment_7_C5.csv と experiment_7_refine_C5.csv (scripts/blocking_reassessment.py の load_exp7 と assess)
% 条件 C5, 系列 NoCancel の ERP 最小点の P_block_arrival_stable の, Base の ERP 最小点の値に対する比
\newcommand{\ExpSevenCFiveNoCancelPRatio}{0.93}

% 出典: results/experiment_7/experiment_7_C5.csv と experiment_7_refine_C5.csv (scripts/blocking_reassessment.py の load_exp7 と assess)
% 条件 C5, 系列 NoCancel の同じ信頼性での改善: P ≤ P_ref (Base の ERP 最小点の P) を
% 満たす点の中の ERP 最小値の, Base の ERP 最小値に対する −(相対差)
\newcommand{\ExpSevenCFiveNoCancelSameRelImprovement}{$-0.99$\%}

% 出典: results/experiment_7/experiment_7_C5.csv と experiment_7_refine_C5.csv (scripts/blocking_reassessment.py の load_exp7 と assess)
% 条件 C5, 系列 P2 only の ERP 最小値 (Base は β を細かくした 121 点を含む)
\newcommand{\ExpSevenCFivePTwoOnlyERPMin}{4.2843}

% 出典: results/experiment_7/experiment_7_C5.csv と experiment_7_refine_C5.csv (scripts/blocking_reassessment.py の load_exp7 と assess)
% 条件 C5, 系列 P2 only の ERP 最小値の, Base の ERP 最小値に対する相対差
\newcommand{\ExpSevenCFivePTwoOnlyERPRel}{$-0.53$\%}

% 出典: results/experiment_7/experiment_7_C5.csv と experiment_7_refine_C5.csv (scripts/blocking_reassessment.py の load_exp7 と assess)
% 条件 C5, 系列 P2 only の ERP 最小点の P_block_arrival_stable の, Base の ERP 最小点の値に対する比
\newcommand{\ExpSevenCFivePTwoOnlyPRatio}{0.95}

% 出典: results/experiment_7/experiment_7_C5.csv と experiment_7_refine_C5.csv (scripts/blocking_reassessment.py の load_exp7 と assess)
% 条件 C5, 系列 P2 only の同じ信頼性での改善: P ≤ P_ref (Base の ERP 最小点の P) を
% 満たす点の中の ERP 最小値の, Base の ERP 最小値に対する −(相対差)
\newcommand{\ExpSevenCFivePTwoOnlySameRelImprovement}{$+0.53$\%}

% 出典: results/experiment_7/experiment_7_C5.csv と experiment_7_refine_C5.csv (scripts/blocking_reassessment.py の load_exp7 と assess)
% 条件 C5, 系列 P1 only の ERP 最小値 (Base は β を細かくした 121 点を含む)
\newcommand{\ExpSevenCFivePOneOnlyERPMin}{4.3069}

% 出典: results/experiment_7/experiment_7_C5.csv と experiment_7_refine_C5.csv (scripts/blocking_reassessment.py の load_exp7 と assess)
% 条件 C5, 系列 P1 only の ERP 最小値の, Base の ERP 最小値に対する相対差
\newcommand{\ExpSevenCFivePOneOnlyERPRel}{$-0.00$\%}

% 出典: results/experiment_7/experiment_7_C5.csv と experiment_7_refine_C5.csv (scripts/blocking_reassessment.py の load_exp7 と assess)
% 条件 C5, 系列 P1 only の ERP 最小点の P_block_arrival_stable の, Base の ERP 最小点の値に対する比
\newcommand{\ExpSevenCFivePOneOnlyPRatio}{1.00}

% 出典: results/experiment_7/experiment_7_C5.csv と experiment_7_refine_C5.csv (scripts/blocking_reassessment.py の load_exp7 と assess)
% 条件 C5, 系列 P1 only の同じ信頼性での改善: P ≤ P_ref (Base の ERP 最小点の P) を
% 満たす点の中の ERP 最小値の, Base の ERP 最小値に対する −(相対差)
\newcommand{\ExpSevenCFivePOneOnlySameRelImprovement}{$+0.00$\%}

% 出典: results/experiment_7/experiment_7_C5.csv と experiment_7_refine_C5.csv (scripts/blocking_reassessment.py の load_exp7 と assess)
% 条件 C5, 系列 P1 only (protect) の ERP 最小値 (Base は β を細かくした 121 点を含む)
\newcommand{\ExpSevenCFivePOneOnlyProtectERPMin}{4.3050}

% 出典: results/experiment_7/experiment_7_C5.csv と experiment_7_refine_C5.csv (scripts/blocking_reassessment.py の load_exp7 と assess)
% 条件 C5, 系列 P1 only (protect) の ERP 最小値の, Base の ERP 最小値に対する相対差
\newcommand{\ExpSevenCFivePOneOnlyProtectERPRel}{$-0.05$\%}

% 出典: results/experiment_7/experiment_7_C5.csv と experiment_7_refine_C5.csv (scripts/blocking_reassessment.py の load_exp7 と assess)
% 条件 C5, 系列 P1 only (protect) の ERP 最小点の P_block_arrival_stable の, Base の ERP 最小点の値に対する比
\newcommand{\ExpSevenCFivePOneOnlyProtectPRatio}{1.00}

% 出典: results/experiment_7/experiment_7_C5.csv と experiment_7_refine_C5.csv (scripts/blocking_reassessment.py の load_exp7 と assess)
% 条件 C5, 系列 P1 only (protect) の同じ信頼性での改善: P ≤ P_ref (Base の ERP 最小点の P) を
% 満たす点の中の ERP 最小値の, Base の ERP 最小値に対する −(相対差)
\newcommand{\ExpSevenCFivePOneOnlyProtectSameRelImprovement}{$+0.05$\%}

% 出典: results/experiment_7/experiment_7_C5.csv と experiment_7_refine_C5.csv (scripts/blocking_reassessment.py の load_exp7 と assess)
% 条件 C5, 系列 P1+P2 の ERP 最小値 (Base は β を細かくした 121 点を含む)
\newcommand{\ExpSevenCFivePOneTwoERPMin}{4.2836}

% 出典: results/experiment_7/experiment_7_C5.csv と experiment_7_refine_C5.csv (scripts/blocking_reassessment.py の load_exp7 と assess)
% 条件 C5, 系列 P1+P2 の ERP 最小値の, Base の ERP 最小値に対する相対差
\newcommand{\ExpSevenCFivePOneTwoERPRel}{$-0.54$\%}

% 出典: results/experiment_7/experiment_7_C5.csv と experiment_7_refine_C5.csv (scripts/blocking_reassessment.py の load_exp7 と assess)
% 条件 C5, 系列 P1+P2 の ERP 最小点の P_block_arrival_stable の, Base の ERP 最小点の値に対する比
\newcommand{\ExpSevenCFivePOneTwoPRatio}{0.95}

% 出典: results/experiment_7/experiment_7_C5.csv と experiment_7_refine_C5.csv (scripts/blocking_reassessment.py の load_exp7 と assess)
% 条件 C5, 系列 P1+P2 の同じ信頼性での改善: P ≤ P_ref (Base の ERP 最小点の P) を
% 満たす点の中の ERP 最小値の, Base の ERP 最小値に対する −(相対差)
\newcommand{\ExpSevenCFivePOneTwoSameRelImprovement}{$+0.54$\%}

% 出典: results/experiment_7/experiment_7_C5.csv と experiment_7_refine_C5.csv (scripts/blocking_reassessment.py の load_exp7 と assess)
% 条件 C5, 系列 P1+P2 (protect) の ERP 最小値 (Base は β を細かくした 121 点を含む)
\newcommand{\ExpSevenCFivePOneTwoProtectERPMin}{4.2833}

% 出典: results/experiment_7/experiment_7_C5.csv と experiment_7_refine_C5.csv (scripts/blocking_reassessment.py の load_exp7 と assess)
% 条件 C5, 系列 P1+P2 (protect) の ERP 最小値の, Base の ERP 最小値に対する相対差
\newcommand{\ExpSevenCFivePOneTwoProtectERPRel}{$-0.55$\%}

% 出典: results/experiment_7/experiment_7_C5.csv と experiment_7_refine_C5.csv (scripts/blocking_reassessment.py の load_exp7 と assess)
% 条件 C5, 系列 P1+P2 (protect) の ERP 最小点の P_block_arrival_stable の, Base の ERP 最小点の値に対する比
\newcommand{\ExpSevenCFivePOneTwoProtectPRatio}{0.95}

% 出典: results/experiment_7/experiment_7_C5.csv と experiment_7_refine_C5.csv (scripts/blocking_reassessment.py の load_exp7 と assess)
% 条件 C5, 系列 P1+P2 (protect) の同じ信頼性での改善: P ≤ P_ref (Base の ERP 最小点の P) を
% 満たす点の中の ERP 最小値の, Base の ERP 最小値に対する −(相対差)
\newcommand{\ExpSevenCFivePOneTwoProtectSameRelImprovement}{$+0.55$\%}

% 出典: results/experiment_7/experiment_7_C5.csv と experiment_7_refine_C5.csv (scripts/blocking_reassessment.py の load_exp7 と assess)
% 条件 C5, Base の ERP 最小点の P_block_arrival_stable (= P_ref)
\newcommand{\ExpSevenCFiveBasePBlock}{1.46e-08}

% 出典: results/experiment_7/experiment_7_C5.csv と experiment_7_refine_C5.csv (scripts/blocking_reassessment.py の load_exp7 と assess)
% 条件 C5, P_block_arrival_stable ≤ 1e-2 を満たす点をもつ系列の数 (Base を含む 7 系列中)
\newcommand{\ExpSevenCFiveEpsTwoFeasibleCount}{7}

% 出典: results/experiment_7/ の CSV (scripts/blocking_reassessment.py の assess)
% 位相を使う 5 系列 (P2 only, P1 only, P1 only (protect), P1+P2, P1+P2 (protect)) と 5 条件で
% ERP 最小値の改善 (−相対差, 対 β を細かくした Base) の最大値
\newcommand{\ExpSevenMaxERPImprovement}{1.19\%}

% 出典: results/experiment_7/ の CSV (scripts/blocking_reassessment.py の assess)
% 位相を使う 5 系列 (P2 only, P1 only, P1 only (protect), P1+P2, P1+P2 (protect)) と 5 条件で
% 同じ信頼性 (P ≤ P_ref) での改善の最大値
\newcommand{\ExpSevenMaxSameRelImprovement}{0.81\%}

% 出典: results/experiment_7/ の CSV (scripts/blocking_reassessment.py の assess)
% NoCancel の ERP 最小値の, Base に対する相対差の最小値 (5 条件すべてで正)
\newcommand{\ExpSevenNoCancelWorseMin}{0.33\%}

% 出典: results/experiment_7/ の CSV (scripts/blocking_reassessment.py の assess)
% 同じく最大値
\newcommand{\ExpSevenNoCancelWorseMax}{8.56\%}

% 出典: results/experiment_{1..4}*.csv と experiment_{1P..4P}_*_{g5.0,g5.0}.csv (scripts/evaluate_predictive_comparison.py の load_experiment_*)
% 旧比較 (保護なしの旧モデル, Base は β を走査の値に固定, Predictive は n_target=10, γ=5) の ERP の改善率
% 405 点中の最大値
\newcommand{\ExpOldBetaFixedGammaFiveMaxERPImprovement}{11.77\%}

% 出典: results/experiment_{1..4}*.csv と experiment_{1P..4P}_*_{g5.0,g5.0}.csv (scripts/evaluate_predictive_comparison.py の load_experiment_*)
% 旧比較 (保護なしの旧モデル, Base は β を走査の値に固定, Predictive は n_target=10, γ=5) の ERP の改善率
% 改善率が 5% を超える点の数 (405 点中)
\newcommand{\ExpOldBetaFixedGammaFiveCountOverFive}{85}

% 出典: results/experiment_{1..4}*.csv と experiment_{1P..4P}_*_{g5.0,g5.0}.csv (scripts/evaluate_predictive_comparison.py の load_experiment_*)
% 旧比較 (保護なしの旧モデル, Base は β を走査の値に固定, Predictive は n_target=10, γ=5) の ERP の改善率
% 405 点の平均
\newcommand{\ExpOldBetaFixedGammaFiveMeanERPImprovement}{$-1.77$\%}

% 出典: results/experiment_{1..4}*.csv と experiment_{1P..4P}_*_{g1.0,gmap}.csv (scripts/evaluate_predictive_comparison.py の load_experiment_*)
% 旧比較 (保護なしの旧モデル, Base は β を走査の値に固定, Predictive は n_target=10, γ を最適化) の ERP の改善率
% 405 点中の最大値
\newcommand{\ExpOldBetaFixedGammaOptMaxERPImprovement}{28.14\%}

% 出典: results/experiment_{1..4}*.csv と experiment_{1P..4P}_*_{g1.0,gmap}.csv (scripts/evaluate_predictive_comparison.py の load_experiment_*)
% 旧比較 (保護なしの旧モデル, Base は β を走査の値に固定, Predictive は n_target=10, γ を最適化) の ERP の改善率
% 改善率が 5% を超える点の数 (405 点中)
\newcommand{\ExpOldBetaFixedGammaOptCountOverFive}{66}

% 出典: results/experiment_{1..4}*.csv と experiment_{1P..4P}_*_{g1.0,gmap}.csv (scripts/evaluate_predictive_comparison.py の load_experiment_*)
% 旧比較 (保護なしの旧モデル, Base は β を走査の値に固定, Predictive は n_target=10, γ を最適化) の ERP の改善率
% 405 点の平均
\newcommand{\ExpOldBetaFixedGammaOptMeanERPImprovement}{$+0.38$\%}

% 出典: results/p1_protect/p1_protect.csv
% 条件 medium, ρ=0.3, α=0.1, β=0.5 で, 保護なし (none, n_target=10, γ=1) の ERP に対する相対差
% 保護の設定 presetup (n_target=10, γ=1)
\newcommand{\ProtectExamplePresetupERPChange}{$-0.71$\%}

% 出典: results/p1_protect/p1_protect.csv
% 条件 medium, ρ=0.3, α=0.1, β=0.5 で, 保護なし (none, n_target=10, γ=1) の ERP に対する相対差
% 保護の設定 both (n_target=10, γ=1)
\newcommand{\ProtectExampleBothERPChange}{$-6.55$\%}

% 出典: results/p1_protect/p1_protect.csv
% 条件 medium, ρ=0.3, α=0.1, β=0.5 で, 保護なし (none, n_target=10, γ=1) の ERP に対する相対差
% NoCancel (never_cancel_setup, n_target=0, γ=1)
\newcommand{\ProtectExampleNoCancelERPChange}{$+12.73$\%}

% 出典: results/p1_protect/p1_protect.csv
% NoCancel の ERP が同じ (条件, α, β) の none (γ=1) より大きい組の数
\newcommand{\ProtectNoCancelWorseCount}{17}

% 出典: results/p1_protect/p1_protect.csv
% 上の組の総数
\newcommand{\ProtectNoCancelTotal}{18}

% 出典: results/experiment_8/experiment_8_<行>.csv (scripts/experiment_8_map.py の summaries)
% 3 系列のいずれかで改善 (−ERP の相対差, 対 β を細かくした Base) が 0.5% を超えるマスの数
\newcommand{\ExpEightCoreCellCount}{13}

% 出典: results/experiment_8/experiment_8_<行>.csv (scripts/experiment_8_map.py の summaries)
% 地図 A・B の重複のないマスの数
\newcommand{\ExpEightCellCount}{35}

% 出典: results/experiment_8/experiment_8_<行>.csv (scripts/experiment_8_map.py の summaries)
% 改善が 0.5% を超えるマスでの, 3 系列のうち最良の改善の最小値
\newcommand{\ExpEightCoreImprovementMin}{0.51\%}

% 出典: results/experiment_8/experiment_8_<行>.csv (scripts/experiment_8_map.py の summaries)
% 改善が 0.5% を超えるマスでの, 3 系列のうち最良の改善の最大値
\newcommand{\ExpEightCoreImprovementMax}{3.79\%}

% 出典: results/experiment_8/experiment_8_<行>.csv (scripts/experiment_8_map.py の summaries)
% 中心の領域 = r1=17.8 の列のうち ρ_B ≤ 0.8 かつ α ≤ 1 のマスと, α=1, r1=3.16 のマス (地図 A・B で重複のないもの)
% 該当: A_rhoB0.6 r1=17.8, A_rhoB0.8 r1=17.8, B_alpha0.005 r1=17.8, B_alpha0.02 r1=17.8, B_alpha1 r1=3.16, B_alpha1 r1=17.8
% マスの数
\newcommand{\ExpEightCenterCellCount}{6}

% 出典: results/experiment_8/experiment_8_<行>.csv (scripts/experiment_8_map.py の summaries)
% 中心の領域 = r1=17.8 の列のうち ρ_B ≤ 0.8 かつ α ≤ 1 のマスと, α=1, r1=3.16 のマス (地図 A・B で重複のないもの)
% 該当: A_rhoB0.6 r1=17.8, A_rhoB0.8 r1=17.8, B_alpha0.005 r1=17.8, B_alpha0.02 r1=17.8, B_alpha1 r1=3.16, B_alpha1 r1=17.8
% 3 系列のうち最良の改善 (−ERP の相対差) の最小値
\newcommand{\ExpEightCenterImprovementMin}{2.24\%}

% 出典: results/experiment_8/experiment_8_<行>.csv (scripts/experiment_8_map.py の summaries)
% 中心の領域 = r1=17.8 の列のうち ρ_B ≤ 0.8 かつ α ≤ 1 のマスと, α=1, r1=3.16 のマス (地図 A・B で重複のないもの)
% 該当: A_rhoB0.6 r1=17.8, A_rhoB0.8 r1=17.8, B_alpha0.005 r1=17.8, B_alpha0.02 r1=17.8, B_alpha1 r1=3.16, B_alpha1 r1=17.8
% 同じく最大値
\newcommand{\ExpEightCenterImprovementMax}{3.79\%}

% 出典: results/experiment_8/experiment_8_<行>.csv (scripts/experiment_8_map.py の summaries)
% 全マス・全系列で改善 (−ERP の相対差) の最大値
\newcommand{\ExpEightMaxImprovement}{3.79\%}

% 出典: results/experiment_8/experiment_8_<行>.csv (scripts/experiment_8_map.py の summaries)
% 改善が最大の系列
\newcommand{\ExpEightMaxImprovementSeries}{P1 only (protect)}

% 出典: results/experiment_8/experiment_8_<行>.csv (scripts/experiment_8_map.py の summaries)
% 改善が最大のマス
\newcommand{\ExpEightMaxImprovementCell}{$\alpha=1$, $r_1=3.16$}

% 出典: results/experiment_8/experiment_8_<行>.csv (scripts/experiment_8_map.py の summaries)
% 系列 P2 only の改善の最大値
\newcommand{\ExpEightPTwoOnlyMaxImprovement}{2.69\%}

% 出典: results/experiment_8/experiment_8_<行>.csv (scripts/experiment_8_map.py の summaries)
% 系列 P2 only の改善が最大のマス
\newcommand{\ExpEightPTwoOnlyMaxImprovementCell}{$\alpha=0.005$, $r_1=17.8$}

% 出典: results/experiment_8/experiment_8_<行>.csv (scripts/experiment_8_map.py の summaries)
% 系列 P2 only の改善が 0.5% を超えるマスの数
\newcommand{\ExpEightPTwoOnlyCoreCellCount}{7}

% 出典: results/experiment_8/experiment_8_<行>.csv (scripts/experiment_8_map.py の summaries)
% 系列 P1 only (protect) の改善の最大値
\newcommand{\ExpEightPOneOnlyProtectMaxImprovement}{3.79\%}

% 出典: results/experiment_8/experiment_8_<行>.csv (scripts/experiment_8_map.py の summaries)
% 系列 P1 only (protect) の改善が最大のマス
\newcommand{\ExpEightPOneOnlyProtectMaxImprovementCell}{$\alpha=1$, $r_1=3.16$}

% 出典: results/experiment_8/experiment_8_<行>.csv (scripts/experiment_8_map.py の summaries)
% 系列 P1 only (protect) の改善が 0.5% を超えるマスの数
\newcommand{\ExpEightPOneOnlyProtectCoreCellCount}{12}

% 出典: results/experiment_8/experiment_8_<行>.csv (scripts/experiment_8_map.py の summaries)
% 系列 P1+P2 (protect) の改善の最大値
\newcommand{\ExpEightPOneTwoProtectMaxImprovement}{2.83\%}

% 出典: results/experiment_8/experiment_8_<行>.csv (scripts/experiment_8_map.py の summaries)
% 系列 P1+P2 (protect) の改善が最大のマス
\newcommand{\ExpEightPOneTwoProtectMaxImprovementCell}{$\alpha=0.005$, $r_1=17.8$}

% 出典: results/experiment_8/experiment_8_<行>.csv (scripts/experiment_8_map.py の summaries)
% 系列 P1+P2 (protect) の改善が 0.5% を超えるマスの数
\newcommand{\ExpEightPOneTwoProtectCoreCellCount}{10}

% 出典: results/experiment_8/experiment_8_<行>.csv (scripts/experiment_8_map.py の summaries)
% r1=0.1 の列の 7 マス × 3 系列での改善の最大値
\newcommand{\ExpEightROneMinMaxImprovement}{0.40\%}

% 出典: results/experiment_8/experiment_8_<行>.csv (scripts/experiment_8_map.py の summaries)
% R = β*(P1 only (protect)) / β*(Base) と P1 only (protect) の改善の相関 (results/experiment_8/mechanism.md の手順 1 と同じ)
% 全マス, Pearson
\newcommand{\ExpEightRPearsonAll}{+0.785}

% 出典: results/experiment_8/experiment_8_<行>.csv (scripts/experiment_8_map.py の summaries)
% R = β*(P1 only (protect)) / β*(Base) と P1 only (protect) の改善の相関 (results/experiment_8/mechanism.md の手順 1 と同じ)
% 全マス, Spearman
\newcommand{\ExpEightRSpearmanAll}{+0.644}

% 出典: results/experiment_8/experiment_8_<行>.csv (scripts/experiment_8_map.py の summaries)
% R = β*(P1 only (protect)) / β*(Base) と P1 only (protect) の改善の相関 (results/experiment_8/mechanism.md の手順 1 と同じ)
% 改善が 0.5% を超えるマス, Pearson
\newcommand{\ExpEightRPearsonCore}{+0.666}

% 出典: results/experiment_8/experiment_8_<行>.csv (scripts/experiment_8_map.py の summaries)
% R = β*(P1 only (protect)) / β*(Base) と P1 only (protect) の改善の相関 (results/experiment_8/mechanism.md の手順 1 と同じ)
% 改善が 0.5% を超えるマス, Spearman
\newcommand{\ExpEightRSpearmanCore}{+0.458}

% 出典: results/experiment_8/experiment_8_<行>.csv (scripts/experiment_8_map.py の summaries)
% R = β*(P1 only (protect)) / β*(Base) と P1 only (protect) の改善の相関 (results/experiment_8/mechanism.md の手順 1 と同じ)
% 改善が 0.5% を超えるマスの数 (P1 only (protect))
\newcommand{\ExpEightRCoreCount}{12}

% 出典: results/experiment_8/experiment_8_<行>.csv (scripts/experiment_8_map.py の summaries)
% R = β*(P1 only (protect)) / β*(Base) と P1 only (protect) の改善の相関 (results/experiment_8/mechanism.md の手順 1 と同じ)
% 改善が 0.5% を超えるマスのうち R > 1 のマスの数
\newcommand{\ExpEightRAboveOneCore}{11}

% 出典: results/experiment_8/experiment_8_<行>.csv (scripts/experiment_8_map.py の summaries)
% P1 only (protect) の ERP 最小点の p1_fire_rate と 1/(1+1/r1) の Pearson の相関係数
\newcommand{\ExpEightFirePearsonRaw}{-0.200}

% 出典: results/experiment_8/experiment_8_<行>.csv (scripts/experiment_8_map.py の summaries)
% 位相 0→1 の遷移 1 回あたりの発動確率 p1_fire_rate/(σϖ₀) と 1/(1+1/r1) の相関
\newcommand{\ExpEightFirePearsonPerOnset}{+0.821}

% 出典: results/experiment_8/experiment_8_<行>.csv (scripts/experiment_8_map.py の summaries)
% 地図 B, r1=17.8 での log10 α と P1 only (protect) の改善の Spearman の順位相関
\newcommand{\ExpEightAlphaSpearmanROneSeventeen}{-1.00}

% 出典: results/experiment_8/experiment_8_<行>.csv (scripts/experiment_8_map.py の summaries)
% 地図 B, r1=100 での log10 α と P1 only (protect) の改善の Spearman の順位相関
\newcommand{\ExpEightAlphaSpearmanROneHundred}{-1.00}

% 出典: results/experiment_8/experiment_8_<行>.csv (scripts/blocking_reassessment.py の assess)
% 同じ信頼性 (P ≤ P_ref) の中の ERP 最小値の, Base の ERP 最小値に対する −(相対差)
% 系列 P2 only の最大値
\newcommand{\ExpEightPTwoOnlySameRelMaxImprovement}{2.22\%}

% 出典: results/experiment_8/experiment_8_<行>.csv (scripts/blocking_reassessment.py の assess)
% 同じ信頼性 (P ≤ P_ref) の中の ERP 最小値の, Base の ERP 最小値に対する −(相対差)
% 系列 P2 only で 0.1% を超えて改善するマスの数
\newcommand{\ExpEightPTwoOnlySameRelImprovedCells}{14}

% 出典: results/experiment_8/ の CSV (assess)
% 系列 P2 only の ERP 最小点の P の, Base の ERP 最小点の P に対する比の最大値
\newcommand{\ExpEightPTwoOnlyPRatioMax}{1.63}

% 出典: results/experiment_8/experiment_8_<行>.csv (scripts/blocking_reassessment.py の assess)
% 同じ信頼性 (P ≤ P_ref) の中の ERP 最小値の, Base の ERP 最小値に対する −(相対差)
% 系列 P1 only (protect) の最大値
\newcommand{\ExpEightPOneOnlyProtectSameRelMaxImprovement}{3.79\%}

% 出典: results/experiment_8/experiment_8_<行>.csv (scripts/blocking_reassessment.py の assess)
% 同じ信頼性 (P ≤ P_ref) の中の ERP 最小値の, Base の ERP 最小値に対する −(相対差)
% 系列 P1 only (protect) で 0.1% を超えて改善するマスの数
\newcommand{\ExpEightPOneOnlyProtectSameRelImprovedCells}{16}

% 出典: results/experiment_8/ の CSV (assess)
% 系列 P1 only (protect) の ERP 最小点の P の, Base の ERP 最小点の P に対する比の最大値
\newcommand{\ExpEightPOneOnlyProtectPRatioMax}{1.82}

% 出典: results/experiment_8/experiment_8_<行>.csv (scripts/blocking_reassessment.py の assess)
% 同じ信頼性 (P ≤ P_ref) の中の ERP 最小値の, Base の ERP 最小値に対する −(相対差)
% 系列 P1+P2 (protect) の最大値
\newcommand{\ExpEightPOneTwoProtectSameRelMaxImprovement}{2.76\%}

% 出典: results/experiment_8/experiment_8_<行>.csv (scripts/blocking_reassessment.py の assess)
% 同じ信頼性 (P ≤ P_ref) の中の ERP 最小値の, Base の ERP 最小値に対する −(相対差)
% 系列 P1+P2 (protect) で 0.1% を超えて改善するマスの数
\newcommand{\ExpEightPOneTwoProtectSameRelImprovedCells}{19}

% 出典: results/experiment_8/ の CSV (assess)
% 系列 P1+P2 (protect) の ERP 最小点の P の, Base の ERP 最小点の P に対する比の最大値
\newcommand{\ExpEightPOneTwoProtectPRatioMax}{2.54}

% 出典: results/experiment_8/ の CSV (assess)
% Base も 3 系列も P ≤ 1e-2 を満たせないマスの数
% 該当: A_rhoB1.12 r1=0.562, A_rhoB1.12 r1=3.16, A_rhoB1.12 r1=17.8, A_rhoB1.12 r1=100
\newcommand{\ExpEightBaseInfeasibleEpsTwoCount}{4}

% 出典: results/experiment_8/blocking_reassessment.md と同じ計算
% ERP だけの判定と同じ信頼性での判定で, 区分 (改善 > 0.1%, 同等, 悪化, 該当なし) が変わった (マス, 系列) の数
\newcommand{\ExpEightCategoryChangedCount}{13}

% 出典: results/experiment_8/mechanism_points.csv
% 実験 7 の C1 と実験 8 で改善が最大の 3 マスの, P1 only (protect) − Base (ERP 最小点どうし)
% 通常期 F=0 のアイドル台数 E[I]_F0 の減少量の最小値
\newcommand{\MechIdleFZeroReductionMin}{0.62}

% 出典: results/experiment_8/mechanism_points.csv
% 実験 7 の C1 と実験 8 で改善が最大の 3 マスの, P1 only (protect) − Base (ERP 最小点どうし)
% 同じく最大値
\newcommand{\MechIdleFZeroReductionMax}{0.91}

% 出典: results/experiment_8/mechanism_points.csv
% 実験 7 の C1 と実験 8 で改善が最大の 3 マスの, P1 only (protect) − Base (ERP 最小点どうし)
% 通常期の Cost_F0 の差の最小値
\newcommand{\MechCostFZeroChangeMin}{$-0.20$}

% 出典: results/experiment_8/mechanism_points.csv
% 実験 7 の C1 と実験 8 で改善が最大の 3 マスの, P1 only (protect) − Base (ERP 最小点どうし)
% 同じく最大値
\newcommand{\MechCostFZeroChangeMax}{$-0.10$}

% 出典: results/experiment_8/mechanism_points.csv
% 実験 7 の C1 と実験 8 で改善が最大の 3 マスの, P1 only (protect) − Base (ERP 最小点どうし)
% バースト期の Cost_F1 の差の最小値
\newcommand{\MechCostFOneChangeMin}{$+0.08$}

% 出典: results/experiment_8/mechanism_points.csv
% 実験 7 の C1 と実験 8 で改善が最大の 3 マスの, P1 only (protect) − Base (ERP 最小点どうし)
% 同じく最大値
\newcommand{\MechCostFOneChangeMax}{$+0.14$}

% 出典: results/experiment_8/mechanism_points.csv
% 実験 7 の C1 と実験 8 で改善が最大の 3 マスの, P1 only (protect) − Base (ERP 最小点どうし)
% Cost の相対差の最小値
\newcommand{\MechCostChangeMin}{$-1.27$\%}

% 出典: results/experiment_8/mechanism_points.csv
% 実験 7 の C1 と実験 8 で改善が最大の 3 マスの, P1 only (protect) − Base (ERP 最小点どうし)
% Cost の相対差の最大値
\newcommand{\MechCostChangeMax}{$-0.38$\%}

% 出典: results/experiment_9/experiment_9_B.csv (しきい値付きの 5 区分: scripts/experiment_9_sensitivity.py の category5;
% |ERP の差| ≤ 0.1%, P の比 0.95〜1.05 を同等. 比較の相手は同じ β* の Base)
% 設定 P1 only (protect) n_target=10 で区分「改善」のマスの数
\newcommand{\ExpNinePOneNTenImprovedCount}{7}

% 出典: results/experiment_9/experiment_9_B.csv (しきい値付きの 5 区分: scripts/experiment_9_sensitivity.py の category5;
% |ERP の差| ≤ 0.1%, P の比 0.95〜1.05 を同等. 比較の相手は同じ β* の Base)
% 設定 P1 only (protect) n_target=10 で区分「同等」のマスの数
\newcommand{\ExpNinePOneNTenEqualCount}{21}

% 出典: results/experiment_9/experiment_9_B.csv (しきい値付きの 5 区分: scripts/experiment_9_sensitivity.py の category5;
% |ERP の差| ≤ 0.1%, P の比 0.95〜1.05 を同等. 比較の相手は同じ β* の Base)
% 設定 P1 only (protect) n_target=10 で区分「棄却のみ改善」のマスの数
\newcommand{\ExpNinePOneNTenBlockingOnlyCount}{3}

% 出典: results/experiment_9/experiment_9_B.csv (しきい値付きの 5 区分: scripts/experiment_9_sensitivity.py の category5;
% |ERP の差| ≤ 0.1%, P の比 0.95〜1.05 を同等. 比較の相手は同じ β* の Base)
% 設定 P1 only (protect) n_target=10 で区分「引き換え」のマスの数
\newcommand{\ExpNinePOneNTenTradeOffCount}{0}

% 出典: results/experiment_9/experiment_9_B.csv (しきい値付きの 5 区分: scripts/experiment_9_sensitivity.py の category5;
% |ERP の差| ≤ 0.1%, P の比 0.95〜1.05 を同等. 比較の相手は同じ β* の Base)
% 設定 P1 only (protect) n_target=10 で区分「悪化」のマスの数
\newcommand{\ExpNinePOneNTenWorseCount}{4}

% 出典: results/experiment_9/experiment_9_B.csv (しきい値付きの 5 区分: scripts/experiment_9_sensitivity.py の category5;
% |ERP の差| ≤ 0.1%, P の比 0.95〜1.05 を同等. 比較の相手は同じ β* の Base)
% 設定 P1 only (protect) n_target=10 で P の比 (対同じ β* の Base) が 1.05 を超えるマスの数 (35 マス中)
\newcommand{\ExpNinePOneNTenPUpCount}{0}

% 出典: results/experiment_9/experiment_9_B.csv (しきい値付きの 5 区分: scripts/experiment_9_sensitivity.py の category5;
% |ERP の差| ≤ 0.1%, P の比 0.95〜1.05 を同等. 比較の相手は同じ β* の Base)
% 設定 P1 only (protect) n_target=10 の, 区分「改善」の中で ERP が最も下がるマスの −(ERP の相対差)
\newcommand{\ExpNinePOneNTenMaxImprovement}{1.04\%}

% 出典: results/experiment_9/experiment_9_B.csv (しきい値付きの 5 区分: scripts/experiment_9_sensitivity.py の category5;
% |ERP の差| ≤ 0.1%, P の比 0.95〜1.05 を同等. 比較の相手は同じ β* の Base)
% 上のマス
\newcommand{\ExpNinePOneNTenMaxImprovementCell}{$\alpha=1$, $r_1=3.16$}

% 出典: results/experiment_9/experiment_9_B.csv (しきい値付きの 5 区分: scripts/experiment_9_sensitivity.py の category5;
% |ERP の差| ≤ 0.1%, P の比 0.95〜1.05 を同等. 比較の相手は同じ β* の Base)
% 設定 P1 only (protect) n_target=10 の, 区分「悪化」「引き換え」の中で ERP の相対差の最大値
\newcommand{\ExpNinePOneNTenMaxWorseERP}{$+2.32$\%}

% 出典: results/experiment_9/experiment_9_B.csv (しきい値付きの 5 区分: scripts/experiment_9_sensitivity.py の category5;
% |ERP の差| ≤ 0.1%, P の比 0.95〜1.05 を同等. 比較の相手は同じ β* の Base)
% 上のマス
\newcommand{\ExpNinePOneNTenMaxWorseERPCell}{$\alpha=10$, $r_1=0.1$}

% 出典: results/experiment_9/experiment_9_B.csv (しきい値付きの 5 区分: scripts/experiment_9_sensitivity.py の category5;
% |ERP の差| ≤ 0.1%, P の比 0.95〜1.05 を同等. 比較の相手は同じ β* の Base)
% 設定 P1 only (protect) n_target=10 の, 区分「悪化」「引き換え」の中で P の比 (対同じ β* の Base) の最大値
\newcommand{\ExpNinePOneNTenMaxPRatio}{0.977}

% 出典: results/experiment_9/experiment_9_B.csv (しきい値付きの 5 区分: scripts/experiment_9_sensitivity.py の category5;
% |ERP の差| ≤ 0.1%, P の比 0.95〜1.05 を同等. 比較の相手は同じ β* の Base)
% 上のマス
\newcommand{\ExpNinePOneNTenMaxPRatioCell}{$\alpha=10$, $r_1=17.8$}

% 出典: results/experiment_9/experiment_9_B.csv (しきい値付きの 5 区分: scripts/experiment_9_sensitivity.py の category5;
% |ERP の差| ≤ 0.1%, P の比 0.95〜1.05 を同等. 比較の相手は同じ β* の Base)
% 設定 P1 only (protect) n_target=20 で区分「改善」のマスの数
\newcommand{\ExpNinePOneNTwentyImprovedCount}{9}

% 出典: results/experiment_9/experiment_9_B.csv (しきい値付きの 5 区分: scripts/experiment_9_sensitivity.py の category5;
% |ERP の差| ≤ 0.1%, P の比 0.95〜1.05 を同等. 比較の相手は同じ β* の Base)
% 設定 P1 only (protect) n_target=20 で区分「同等」のマスの数
\newcommand{\ExpNinePOneNTwentyEqualCount}{3}

% 出典: results/experiment_9/experiment_9_B.csv (しきい値付きの 5 区分: scripts/experiment_9_sensitivity.py の category5;
% |ERP の差| ≤ 0.1%, P の比 0.95〜1.05 を同等. 比較の相手は同じ β* の Base)
% 設定 P1 only (protect) n_target=20 で区分「棄却のみ改善」のマスの数
\newcommand{\ExpNinePOneNTwentyBlockingOnlyCount}{0}

% 出典: results/experiment_9/experiment_9_B.csv (しきい値付きの 5 区分: scripts/experiment_9_sensitivity.py の category5;
% |ERP の差| ≤ 0.1%, P の比 0.95〜1.05 を同等. 比較の相手は同じ β* の Base)
% 設定 P1 only (protect) n_target=20 で区分「引き換え」のマスの数
\newcommand{\ExpNinePOneNTwentyTradeOffCount}{23}

% 出典: results/experiment_9/experiment_9_B.csv (しきい値付きの 5 区分: scripts/experiment_9_sensitivity.py の category5;
% |ERP の差| ≤ 0.1%, P の比 0.95〜1.05 を同等. 比較の相手は同じ β* の Base)
% 設定 P1 only (protect) n_target=20 で区分「悪化」のマスの数
\newcommand{\ExpNinePOneNTwentyWorseCount}{0}

% 出典: results/experiment_9/experiment_9_B.csv (しきい値付きの 5 区分: scripts/experiment_9_sensitivity.py の category5;
% |ERP の差| ≤ 0.1%, P の比 0.95〜1.05 を同等. 比較の相手は同じ β* の Base)
% 設定 P1 only (protect) n_target=20 で P の比 (対同じ β* の Base) が 1.05 を超えるマスの数 (35 マス中)
\newcommand{\ExpNinePOneNTwentyPUpCount}{0}

% 出典: results/experiment_9/experiment_9_B.csv (しきい値付きの 5 区分: scripts/experiment_9_sensitivity.py の category5;
% |ERP の差| ≤ 0.1%, P の比 0.95〜1.05 を同等. 比較の相手は同じ β* の Base)
% 設定 P1 only (protect) n_target=20 の, 区分「改善」の中で ERP が最も下がるマスの −(ERP の相対差)
\newcommand{\ExpNinePOneNTwentyMaxImprovement}{2.32\%}

% 出典: results/experiment_9/experiment_9_B.csv (しきい値付きの 5 区分: scripts/experiment_9_sensitivity.py の category5;
% |ERP の差| ≤ 0.1%, P の比 0.95〜1.05 を同等. 比較の相手は同じ β* の Base)
% 上のマス
\newcommand{\ExpNinePOneNTwentyMaxImprovementCell}{$\rho_B=0.8$, $r_1=17.8$}

% 出典: results/experiment_9/experiment_9_B.csv (しきい値付きの 5 区分: scripts/experiment_9_sensitivity.py の category5;
% |ERP の差| ≤ 0.1%, P の比 0.95〜1.05 を同等. 比較の相手は同じ β* の Base)
% 設定 P1 only (protect) n_target=20 の, 区分「悪化」「引き換え」の中で ERP の相対差の最大値
\newcommand{\ExpNinePOneNTwentyMaxWorseERP}{$+53.93$\%}

% 出典: results/experiment_9/experiment_9_B.csv (しきい値付きの 5 区分: scripts/experiment_9_sensitivity.py の category5;
% |ERP の差| ≤ 0.1%, P の比 0.95〜1.05 を同等. 比較の相手は同じ β* の Base)
% 上のマス
\newcommand{\ExpNinePOneNTwentyMaxWorseERPCell}{$\alpha=10$, $r_1=0.1$}

% 出典: results/experiment_9/experiment_9_B.csv (しきい値付きの 5 区分: scripts/experiment_9_sensitivity.py の category5;
% |ERP の差| ≤ 0.1%, P の比 0.95〜1.05 を同等. 比較の相手は同じ β* の Base)
% 設定 P1 only (protect) n_target=20 の, 区分「悪化」「引き換え」の中で P の比 (対同じ β* の Base) の最大値
\newcommand{\ExpNinePOneNTwentyMaxPRatio}{0.794}

% 出典: results/experiment_9/experiment_9_B.csv (しきい値付きの 5 区分: scripts/experiment_9_sensitivity.py の category5;
% |ERP の差| ≤ 0.1%, P の比 0.95〜1.05 を同等. 比較の相手は同じ β* の Base)
% 上のマス
\newcommand{\ExpNinePOneNTwentyMaxPRatioCell}{$\alpha=10$, $r_1=100$}

% 出典: results/experiment_9/experiment_9_B.csv (しきい値付きの 5 区分: scripts/experiment_9_sensitivity.py の category5;
% |ERP の差| ≤ 0.1%, P の比 0.95〜1.05 を同等. 比較の相手は同じ β* の Base)
% 設定 P2 only γ=3 で区分「改善」のマスの数
\newcommand{\ExpNinePTwoGammaThreeImprovedCount}{3}

% 出典: results/experiment_9/experiment_9_B.csv (しきい値付きの 5 区分: scripts/experiment_9_sensitivity.py の category5;
% |ERP の差| ≤ 0.1%, P の比 0.95〜1.05 を同等. 比較の相手は同じ β* の Base)
% 設定 P2 only γ=3 で区分「同等」のマスの数
\newcommand{\ExpNinePTwoGammaThreeEqualCount}{3}

% 出典: results/experiment_9/experiment_9_B.csv (しきい値付きの 5 区分: scripts/experiment_9_sensitivity.py の category5;
% |ERP の差| ≤ 0.1%, P の比 0.95〜1.05 を同等. 比較の相手は同じ β* の Base)
% 設定 P2 only γ=3 で区分「棄却のみ改善」のマスの数
\newcommand{\ExpNinePTwoGammaThreeBlockingOnlyCount}{0}

% 出典: results/experiment_9/experiment_9_B.csv (しきい値付きの 5 区分: scripts/experiment_9_sensitivity.py の category5;
% |ERP の差| ≤ 0.1%, P の比 0.95〜1.05 を同等. 比較の相手は同じ β* の Base)
% 設定 P2 only γ=3 で区分「引き換え」のマスの数
\newcommand{\ExpNinePTwoGammaThreeTradeOffCount}{6}

% 出典: results/experiment_9/experiment_9_B.csv (しきい値付きの 5 区分: scripts/experiment_9_sensitivity.py の category5;
% |ERP の差| ≤ 0.1%, P の比 0.95〜1.05 を同等. 比較の相手は同じ β* の Base)
% 設定 P2 only γ=3 で区分「悪化」のマスの数
\newcommand{\ExpNinePTwoGammaThreeWorseCount}{23}

% 出典: results/experiment_9/experiment_9_B.csv (しきい値付きの 5 区分: scripts/experiment_9_sensitivity.py の category5;
% |ERP の差| ≤ 0.1%, P の比 0.95〜1.05 を同等. 比較の相手は同じ β* の Base)
% 設定 P2 only γ=3 で P の比 (対同じ β* の Base) が 1.05 を超えるマスの数 (35 マス中)
\newcommand{\ExpNinePTwoGammaThreePUpCount}{27}

% 出典: results/experiment_9/experiment_9_B.csv (しきい値付きの 5 区分: scripts/experiment_9_sensitivity.py の category5;
% |ERP の差| ≤ 0.1%, P の比 0.95〜1.05 を同等. 比較の相手は同じ β* の Base)
% 設定 P2 only γ=3 の, 区分「改善」の中で ERP が最も下がるマスの −(ERP の相対差)
\newcommand{\ExpNinePTwoGammaThreeMaxImprovement}{0.18\%}

% 出典: results/experiment_9/experiment_9_B.csv (しきい値付きの 5 区分: scripts/experiment_9_sensitivity.py の category5;
% |ERP の差| ≤ 0.1%, P の比 0.95〜1.05 を同等. 比較の相手は同じ β* の Base)
% 上のマス
\newcommand{\ExpNinePTwoGammaThreeMaxImprovementCell}{$\alpha=10$, $r_1=17.8$}

% 出典: results/experiment_9/experiment_9_B.csv (しきい値付きの 5 区分: scripts/experiment_9_sensitivity.py の category5;
% |ERP の差| ≤ 0.1%, P の比 0.95〜1.05 を同等. 比較の相手は同じ β* の Base)
% 設定 P2 only γ=3 の, 区分「悪化」「引き換え」の中で ERP の相対差の最大値
\newcommand{\ExpNinePTwoGammaThreeMaxWorseERP}{$+2.71$\%}

% 出典: results/experiment_9/experiment_9_B.csv (しきい値付きの 5 区分: scripts/experiment_9_sensitivity.py の category5;
% |ERP の差| ≤ 0.1%, P の比 0.95〜1.05 を同等. 比較の相手は同じ β* の Base)
% 上のマス
\newcommand{\ExpNinePTwoGammaThreeMaxWorseERPCell}{$\rho_B=0.6$, $r_1=0.562$}

% 出典: results/experiment_9/experiment_9_B.csv (しきい値付きの 5 区分: scripts/experiment_9_sensitivity.py の category5;
% |ERP の差| ≤ 0.1%, P の比 0.95〜1.05 を同等. 比較の相手は同じ β* の Base)
% 設定 P2 only γ=3 の, 区分「悪化」「引き換え」の中で P の比 (対同じ β* の Base) の最大値
\newcommand{\ExpNinePTwoGammaThreeMaxPRatio}{11.528}

% 出典: results/experiment_9/experiment_9_B.csv (しきい値付きの 5 区分: scripts/experiment_9_sensitivity.py の category5;
% |ERP の差| ≤ 0.1%, P の比 0.95〜1.05 を同等. 比較の相手は同じ β* の Base)
% 上のマス
\newcommand{\ExpNinePTwoGammaThreeMaxPRatioCell}{$\rho_B=0.6$, $r_1=0.562$}

% 出典: results/experiment_9/experiment_9_B.csv (しきい値付きの 5 区分: scripts/experiment_9_sensitivity.py の category5;
% |ERP の差| ≤ 0.1%, P の比 0.95〜1.05 を同等. 比較の相手は同じ β* の Base)
% 設定 P2 only γ=10 で区分「改善」のマスの数
\newcommand{\ExpNinePTwoGammaTenImprovedCount}{2}

% 出典: results/experiment_9/experiment_9_B.csv (しきい値付きの 5 区分: scripts/experiment_9_sensitivity.py の category5;
% |ERP の差| ≤ 0.1%, P の比 0.95〜1.05 を同等. 比較の相手は同じ β* の Base)
% 設定 P2 only γ=10 で区分「同等」のマスの数
\newcommand{\ExpNinePTwoGammaTenEqualCount}{2}

% 出典: results/experiment_9/experiment_9_B.csv (しきい値付きの 5 区分: scripts/experiment_9_sensitivity.py の category5;
% |ERP の差| ≤ 0.1%, P の比 0.95〜1.05 を同等. 比較の相手は同じ β* の Base)
% 設定 P2 only γ=10 で区分「棄却のみ改善」のマスの数
\newcommand{\ExpNinePTwoGammaTenBlockingOnlyCount}{0}

% 出典: results/experiment_9/experiment_9_B.csv (しきい値付きの 5 区分: scripts/experiment_9_sensitivity.py の category5;
% |ERP の差| ≤ 0.1%, P の比 0.95〜1.05 を同等. 比較の相手は同じ β* の Base)
% 設定 P2 only γ=10 で区分「引き換え」のマスの数
\newcommand{\ExpNinePTwoGammaTenTradeOffCount}{0}

% 出典: results/experiment_9/experiment_9_B.csv (しきい値付きの 5 区分: scripts/experiment_9_sensitivity.py の category5;
% |ERP の差| ≤ 0.1%, P の比 0.95〜1.05 を同等. 比較の相手は同じ β* の Base)
% 設定 P2 only γ=10 で区分「悪化」のマスの数
\newcommand{\ExpNinePTwoGammaTenWorseCount}{31}

% 出典: results/experiment_9/experiment_9_B.csv (しきい値付きの 5 区分: scripts/experiment_9_sensitivity.py の category5;
% |ERP の差| ≤ 0.1%, P の比 0.95〜1.05 を同等. 比較の相手は同じ β* の Base)
% 設定 P2 only γ=10 で P の比 (対同じ β* の Base) が 1.05 を超えるマスの数 (35 マス中)
\newcommand{\ExpNinePTwoGammaTenPUpCount}{28}

% 出典: results/experiment_9/experiment_9_B.csv (しきい値付きの 5 区分: scripts/experiment_9_sensitivity.py の category5;
% |ERP の差| ≤ 0.1%, P の比 0.95〜1.05 を同等. 比較の相手は同じ β* の Base)
% 設定 P2 only γ=10 の, 区分「改善」の中で ERP が最も下がるマスの −(ERP の相対差)
\newcommand{\ExpNinePTwoGammaTenMaxImprovement}{0.21\%}

% 出典: results/experiment_9/experiment_9_B.csv (しきい値付きの 5 区分: scripts/experiment_9_sensitivity.py の category5;
% |ERP の差| ≤ 0.1%, P の比 0.95〜1.05 を同等. 比較の相手は同じ β* の Base)
% 上のマス
\newcommand{\ExpNinePTwoGammaTenMaxImprovementCell}{$\alpha=10$, $r_1=17.8$}

% 出典: results/experiment_9/experiment_9_B.csv (しきい値付きの 5 区分: scripts/experiment_9_sensitivity.py の category5;
% |ERP の差| ≤ 0.1%, P の比 0.95〜1.05 を同等. 比較の相手は同じ β* の Base)
% 設定 P2 only γ=10 の, 区分「悪化」「引き換え」の中で ERP の相対差の最大値
\newcommand{\ExpNinePTwoGammaTenMaxWorseERP}{$+17.11$\%}

% 出典: results/experiment_9/experiment_9_B.csv (しきい値付きの 5 区分: scripts/experiment_9_sensitivity.py の category5;
% |ERP の差| ≤ 0.1%, P の比 0.95〜1.05 を同等. 比較の相手は同じ β* の Base)
% 上のマス
\newcommand{\ExpNinePTwoGammaTenMaxWorseERPCell}{$\rho_B=0.6$, $r_1=0.562$}

% 出典: results/experiment_9/experiment_9_B.csv (しきい値付きの 5 区分: scripts/experiment_9_sensitivity.py の category5;
% |ERP の差| ≤ 0.1%, P の比 0.95〜1.05 を同等. 比較の相手は同じ β* の Base)
% 設定 P2 only γ=10 の, 区分「悪化」「引き換え」の中で P の比 (対同じ β* の Base) の最大値
\newcommand{\ExpNinePTwoGammaTenMaxPRatio}{129.275}

% 出典: results/experiment_9/experiment_9_B.csv (しきい値付きの 5 区分: scripts/experiment_9_sensitivity.py の category5;
% |ERP の差| ≤ 0.1%, P の比 0.95〜1.05 を同等. 比較の相手は同じ β* の Base)
% 上のマス
\newcommand{\ExpNinePTwoGammaTenMaxPRatioCell}{$\rho_B=0.6$, $r_1=0.562$}

% 出典: results/experiment_9/experiment_9_A.csv
% 基準点の β を 2 倍にしたとき P の比 (対 P_ref) が 1.05 を超える組の数
\newcommand{\ExpNineBetaTwoOverCount}{8}

% 出典: results/experiment_9/experiment_9_A.csv
% 基準点の組の数
\newcommand{\ExpNineBetaTwoTotal}{12}

% 出典: results/experiment_9/experiment_9_A.csv
% 上の組の P の比の最大値
\newcommand{\ExpNineBetaTwoMaxPRatio}{2.951}

% 出典: results_gth/benchmark.md (2 コア, numba JIT のウォームアップ後の 3 回中の最良値)
% ベースモデル標準設定の状態数 N
\newcommand{\BenchBaseN}{8{,}442}

% 出典: results_gth/benchmark.md (2 コア, numba JIT のウォームアップ後の 3 回中の最良値)
% ベースモデル標準設定の下帯幅 p
\newcommand{\BenchBaseP}{210}

% 出典: results_gth/benchmark.md (2 コア, numba JIT のウォームアップ後の 3 回中の最良値)
% ベースモデル標準設定の上帯幅 q
\newcommand{\BenchBaseQ}{42}

% 出典: results_gth/benchmark.md (2 コア, numba JIT のウォームアップ後の 3 回中の最良値)
% ベースモデル標準設定の GTH の計算時間 [ms]
\newcommand{\BenchBaseGTHms}{102.5}

% 出典: results_gth/benchmark.md (2 コア, numba JIT のウォームアップ後の 3 回中の最良値)
% Predictive 標準設定の名目の状態数
\newcommand{\BenchPredNominalN}{92{,}862}

% 出典: results_gth/benchmark.md (2 コア, numba JIT のウォームアップ後の 3 回中の最良値)
% Predictive 標準設定の到達可能状態に縮約した後の状態数
\newcommand{\BenchPredReducedN}{12{,}292}

% 出典: results_gth/benchmark.md (2 コア, numba JIT のウォームアップ後の 3 回中の最良値)
% Predictive 標準設定の縮約後の下帯幅 p
\newcommand{\BenchPredP}{760}

% 出典: results_gth/benchmark.md (2 コア, numba JIT のウォームアップ後の 3 回中の最良値)
% Predictive 標準設定の縮約後の上帯幅 q
\newcommand{\BenchPredQ}{152}

% 出典: results_gth/benchmark.md (2 コア, numba JIT のウォームアップ後の 3 回中の最良値)
% Predictive 標準設定の GTH の計算時間 [s]
\newcommand{\BenchPredGTHs}{0.97}

% 出典: results_gth/benchmark.md (2 コア, numba JIT のウォームアップ後の 3 回中の最良値)
% Predictive K=1000 の縮約後の状態数
\newcommand{\BenchPredKThousandReducedN}{45{,}892}

% 出典: results_gth/benchmark.md (2 コア, numba JIT のウォームアップ後の 3 回中の最良値)
% Predictive K=1000 の GTH の計算時間 [s]
\newcommand{\BenchPredKThousandGTHs}{2.33}
